\subsection{Change of rings; Macaulay duality}

PROPOSITION 5.--- Let \(\rho: A \to B\) be a local homomorphism of Noetherian local rings, such that the residue field extension \(\kappa_A \to \kappa_B\) induced by \(\rho\) has finite degree. Let \(I_A\) be a Matlis A-module.

\begin{enumerate}
\def\labelenumi{\alph{enumi})}
\item
  Let \(I_{B}\) be the sub-B-module of \(\mathrm{Hom}_{\mathrm{A}}(\mathrm{B},\mathrm{I}_{\mathrm{A}})\) consisting of the A-homomorphisms from B into \(I_{A}\) whose kernel contains a power of \(m_{B}\). Then \(I_{B}\) is a Matlis B-module.
\item
  Let M be a B-module. The canonical map
\end{enumerate}

\[
\alpha : \operatorname{Hom} _ {\mathrm{B}} (\mathrm{M}, \operatorname{Hom} _ {\mathrm{A}} (\mathrm{B}, \mathrm{I} _ {\mathrm{A}})) \longrightarrow \operatorname{Hom} _ {\mathrm{A}} (\mathrm{M}, \mathrm{I} _ {\mathrm{A}})
\]

defined by \(\alpha(u)(m)=u(m)(1)\) induces a B-isomorphism of \(\mathrm{D}_{\mathrm{B}}(\mathrm{M})=\mathrm{Hom}_{\mathrm{B}}(\mathrm{M},\mathrm{I}_{\mathrm{B}})\) on the sub-B-module \(\mathrm{Hom}_{\mathrm{A}}^{\mathrm{cont}}(\mathrm{M},\mathrm{I}_{\mathrm{A}})\) of \(\mathrm{D}_{\mathrm{A}}(\mathrm{M})=\mathrm{Hom}_{\mathrm{A}}(\mathrm{M},\mathrm{I}_{\mathrm{A}})\) formed by the applications \(f:\mathrm{M}\to\mathrm{I}_{\mathrm{A}}\) such that for every element \(m\) of M, there exists an integer \(n\geqslant0\) such that \(f(\mathfrak{m}_{\mathrm{B}}^{n}m)=0\).

The above condition on \(f\) means that \(f\) is continuous when one equips \(I_A\) with the discrete topology and M with the finest topology inducing on each finitely generated submodule the topology \(\mathfrak{m}_B\)-adic, which justifies the notation. Likewise, the condition \(g \in I_B\) means that \(g\) is continuous when one equips \(I_A\) with the discrete topology and B with the topology \(\mathfrak{m}_B\) -adic.

Let us prove a). For every finite-length B-module M, denote by T(M) the B-module \(\mathrm{Hom}_{\mathrm{A}}(\mathrm{M},\mathrm{I}_{\mathrm{A}})\). For every B-linear map \(f:M\to N\) between B-modules of finite length, denote by \(\mathrm{T}(f):\mathrm{T}(\mathrm{N})\to\mathrm{T}(\mathrm{M})\) the B-linear map \(\mathrm{Hom}_{\mathrm{A}}(f,1_{\mathrm{I}_{\mathrm{A}}})\). Verification of conditions FD 1) through FD 4) of no. 5 is immediate. Moreover, for every B-module of finite length N, we have \(\mathrm{long}_{\mathrm{A}}(\mathrm{N})=\mathrm{long}_{\mathrm{B}}(\mathrm{N})\left[\kappa_{\mathrm{B}}:\kappa_{\mathrm{A}}\right]\); since we have \(\mathrm{long}_{\mathrm{A}}(\mathrm{T}(\mathrm{M}))=\mathrm{long}_{\mathrm{A}}(\mathrm{M})\) (no. 3, th. 2), we deduce \(\mathrm{long}_{\mathrm{B}}(\mathrm{T}(\mathrm{M}))=\mathrm{long}_{\mathrm{B}}(\mathrm{M})\), which implies FD 5). One may therefore apply Theorem 3 of no. 5; one has

\[
\mathrm{T} (\mathrm{B} / \mathfrak {m} _ {\mathrm{B}} ^ {n}) = \operatorname{Hom} _ {\mathrm{A}} (\mathrm{B} / \mathfrak {m} _ {\mathrm{B}} ^ {n}, \mathrm{I} _ {\mathrm{A}}),
\]

so that the Matlis B-module \(\varinjlim\mathrm{T}(\mathrm{B}/\mathfrak{m}_{\mathrm{B}}^{n})\) is identified with the sub-B-module \(\mathrm{I}_{\mathrm{B}}\) of \(\operatorname{Hom}_{\mathrm{A}}(\mathrm{B},\mathrm{I}_{\mathrm{A}})\), which proves a).

Let us prove b). The map \(\alpha\) is the inverse of the canonical isomorphism

\[
\beta : \operatorname{Hom} _ {\mathrm{A}} (\mathrm{M}, \mathrm{I} _ {\mathrm{A}}) \longrightarrow \operatorname{Hom} _ {\mathrm{B}} (\mathrm{M}, \operatorname{Hom} _ {\mathrm{A}} (\mathrm{B}, \mathrm{I} _ {\mathrm{A}}))
\]

which associates to \(v \in \mathrm{Hom}_{\mathrm{A}}(\mathrm{M}, \mathrm{I}_{\mathrm{A}})\) the map \(v'\) from M to \(\mathrm{Hom}_{\mathrm{A}}(\mathrm{B}, \mathrm{I}_{\mathrm{A}})\) such that \(v'(m)(b) = v(bm)\) (A, II, p. 74, prop. 1). For \(v'\) to take its values in \(\mathrm{I}_{\mathrm{B}}\), it is necessary and sufficient that \(v\) belong to \(\mathrm{Hom}_{\mathrm{A}}^{cont}(\mathrm{M}, \mathrm{I}_{\mathrm{A}})\) , whence b).

COROLLARY.--- a) If the A-algebra B is finite, the B-module \(\mathrm{I_B} = \mathrm{Hom}_{\mathrm{A}}(\mathrm{B},\mathrm{I}_{\mathrm{A}})\) is a Matlis B-module.

\begin{enumerate}
\def\labelenumi{\alph{enumi})}
\setcounter{enumi}{1}
\tightlist
\item
  If the B-module M is Artinian, the map \(\alpha\) is a B-isomorphism from \(D_{B}(M)\) onto \(D_{A}(M)\).
\end{enumerate}

If the A-algebra B is finite, so is the \(\kappa_{A}\) -algebra \(B/\mathfrak{m}_{A}B\), which implies that \(\mathfrak{m}_{A}B\) is a defining ideal of B (VIII, § 3, n° 2, lemma 2). Since every element of the Matlis module \(I_{A}\) is annihilated by a power of \(\mathfrak{m}_{A}\), every element of \(\mathrm{Hom}_{\mathrm{A}}(\mathrm{B}, I_{\mathrm{A}})\) is annihilated by a power of \(\mathfrak{m}_{\mathrm{B}}\) , whence a). Assertion b) follows from the fact that every element of an Artinian module is annihilated by a power of the maximal ideal (n° 3, lemma 4).

Prop. 5 applies in particular when A is a field \(k\), in which case one may take \(\mathrm{I}_{\mathrm{A}} = k\) , hence \(\mathrm{D}_{k}(\mathrm{M}) = \mathrm{Hom}_{k}^{cont}(\mathrm{M}, k)\) (“Macaulay duality”). Note that the hypothesis \([\kappa_{\mathrm{B}} : k] < +\infty\) is in particular satisfied when the \(k\) -algebra B is the local ring at a maximal ideal of a \(k\) -algebra of finite type (A, VIII, App. 3, cor. 1).

More particularly, consider a \(k\)-algebra of finite type S, graded of type N, such that \(S_0\) is a field, a finite-degree extension of \(k\). One may apply prop. 5 to the local ring \(S'\) of S at the maximal ideal \(S_+ = \bigoplus_{n > 0} S_n\) or, equivalently, to its completion \(\widehat{S} = \prod_{n \geqslant 0} S_n\) (III, § 1, n° 3, lemma 2 and § 2, n° 12, example 1). The S-module \(I_{\widehat{S}} = \text{Hom}_k^{cont}(\widehat{S}, k)\) then identifies with

\[
\mathrm{S} ^ {* \mathrm{gr}} = \bigoplus_ {n \geqslant 0} \operatorname{Hom} _ {k} (\mathrm{S} _ {n}, k) ;
\]

for \(s \in S\) and \(u \in S^{*gr}\), the element \(su\) of \(S^{*gr}\) is the interior product \(s \lrcorner u\) (A, III, p. 156 and p. 157). Take for example \(S = k[T_1, \ldots, T_d]\), whence \(\widehat{S} = k[[T_1, \ldots, T_d]]\). Let \((u_\alpha)_{\alpha \in \mathbf{N}^d}\) denote the basis of the \(k\) -vector space \(S^{*gr}\) dual to the basis \((T^\alpha)_{\alpha \in N^d}\) of S. The S-module structure of \(S^{*gr}\) is then described by the formulas (A, III, p. 167)

\[
\begin{array}{l l} \mathbf {T} ^ {\beta} u _ {\alpha} = u _ {\alpha - \beta} & \text { if } \alpha \geqslant \beta , \\ \mathbf {T} ^ {\beta} u _ {\alpha} = 0 & \text { otherwise. } \end{array}
\]
% semantic-fragments: legacy-input-seam
\relax{}
